% ENDLEAF-R49 and ENDLEAF-R39-KEEP.
% Kernel boxes only. This file is not a package, so packageSetHash
% does not change.
% The template owns A4 page geometry. A sysml canvas wider than the
% portrait text line at the declared body size is landscape A4, and
% type stays at that declared size. 7 pt is an absolute floor, not a
% fit target. Do not scale type down to force portrait. Do not shrink
% the coordinate unit (that leaves names outside their boxes). If the
% canvas still will not fit landscape A4 at the declared size, this
% is an error. Portrait sheets that already fit stay portrait.
% Caller \pdfpagewidth in the body is not the switch: a minipage
% group would restore portrait at shipout.
% Plain tikzpictures are not scaled. With no job orientation a plain
% picture wider than the portrait line turns its page landscape
% (ENDLEAF-R69-AUTO, ENDLEAF_AUTO_LANDSCAPE); a picture still wider than
% its line logs ENDLEAF_WIDE_PICTURE (ENDLEAF-R69-WIDE-WARN).
% ENDLEAF-R59-FLOAT. Inside a figure, a canvas that needs landscape makes
% that float page-only: it is never set here, top or bottom on a portrait
% page and never shares a float page. When it reaches the head of the
% deferred list it ships alone on a landscape page; the next page is
% portrait. Text after the figure keeps flowing: no forced page break.
% ENDLEAF-R59-OVERHANG. A shipped canvas still wider than its line logs
% ENDLEAF_OVERHANG; the worker turns that and ENDLEAF_WIDE_PICTURE into a
% layout_overhang job warning.
% ENDLEAF-R56-FIT. The class paper option sets the size (A4, Letter or
% A3); \EndleafPaper names it in errors. With no job orientation
% (\EndleafOrientation undefined) the choice above is unchanged. With a
% job orientation the page stays as the class set it, and a canvas that
% does not fit that page is an error, never a silent turn.
\NeedsTeXFormat{LaTeX2e}
\makeatletter
\newbox\elfit@picbox
\newbox\elfit@plainbox
\newdimen\elfit@availht
\newdimen\elfit@nath
\newdimen\elfit@portline
\newdimen\elfit@porttw
\newdimen\elfit@portth
\newdimen\elfit@portpw
\newdimen\elfit@portph
\newdimen\elfit@portoddm
\newdimen\elfit@portevenm
\newdimen\elfit@porttopm
\newdimen\elfit@landpw
\newdimen\elfit@landph
\newdimen\elfit@landtw
\newdimen\elfit@landth
\newdimen\elfit@landline
\newif\ifelfit@needland
\newif\ifelfit@thisland
\newif\ifelfit@landscaped
\newif\ifelfit@insidefloat
\newif\ifelfit@locked
\newif\ifelfit@shipland
% ENDLEAF-R80-MARGINS. Pinned landscape margins; endleaf-page.tex defines
% them. Same numbers here when this file runs without it.
\ifdefined\EndleafLandL\else
  \newdimen\EndleafLandL
  \newdimen\EndleafLandR
  \newdimen\EndleafLandT
  \newdimen\EndleafLandB
  \EndleafLandL=20mm
  \EndleafLandR=20mm
  \EndleafLandT=20mm
  \EndleafLandB=22mm
\fi
\ifdefined\EndleafPaper
  \edef\elfit@papername{\EndleafPaper}
\else
  \def\elfit@papername{A4}
\fi
\ifdefined\EndleafOrientation
  \elfit@lockedtrue
  \edef\elfit@pagename{\EndleafOrientation}
\else
  \def\elfit@pagename{portrait}
\fi
\def\elfit@uniform{1}
\def\elfit@unit{1}
\def\elfit@shown{}
\def\elfit@page{portrait}
\let\elfit@tikzpicture\tikzpicture
\let\elfit@endtikzpicture\endtikzpicture
\let\elfit@figure\figure
\let\elfit@endfigure\endfigure
\let\elfit@table\table
\let\elfit@endtable\endtable

% Snapshot portrait article metrics. ENDLEAF-R80-MARGINS: the landscape
% text block is the turned paper minus \EndleafLandL/R/T/B (A4: 257 x
% 168 mm), the same block a named landscape page gets.
\AtBeginDocument{%
  \elfit@portline=\linewidth
  \elfit@porttw=\textwidth
  \elfit@portth=\textheight
  \elfit@portpw=\paperwidth
  \elfit@portph=\paperheight
  \elfit@portoddm=\oddsidemargin
  \elfit@portevenm=\evensidemargin
  \elfit@porttopm=\topmargin
  \elfit@landpw=\paperheight
  \elfit@landph=\paperwidth
  \elfit@landtw=\dimexpr\elfit@landpw-\EndleafLandL-\EndleafLandR\relax
  \elfit@landth=\dimexpr\elfit@landph-\EndleafLandT-\EndleafLandB\relax
  \ifdim\elfit@landth<80pt
    \elfit@landth=80pt
  \fi
  \elfit@landline=\elfit@landtw
}

\def\elfit@applypapersize{%
  \ifdefined\pdfpagewidth
    \global\pdfpagewidth=\paperwidth
    \global\pdfpageheight=\paperheight
  \fi
  \ifdefined\pagewidth
    \global\pagewidth=\paperwidth
    \global\pageheight=\paperheight
  \fi
}

\def\elfit@setlandscape{%
  \global\paperwidth=\elfit@landpw
  \global\paperheight=\elfit@landph
  \global\oddsidemargin=\dimexpr\EndleafLandL-1in\relax
  \global\evensidemargin=\dimexpr\EndleafLandL-1in\relax
  \global\topmargin=\dimexpr\EndleafLandT-1in-\headheight-\headsep\relax
  \global\textwidth=\elfit@landtw
  \global\textheight=\elfit@landth
  \global\hsize=\elfit@landtw
  \global\linewidth=\elfit@landtw
  \global\columnwidth=\elfit@landtw
  \global\vsize=\elfit@landth
  \global\@colht=\elfit@landth
  \global\@colroom=\elfit@landth
  \elfit@applypapersize
  \gdef\elfit@page{landscape}%
  \global\elfit@landscapedtrue
}

\def\elfit@setportrait{%
  \global\paperwidth=\elfit@portpw
  \global\paperheight=\elfit@portph
  \global\oddsidemargin=\elfit@portoddm
  \global\evensidemargin=\elfit@portevenm
  \global\topmargin=\elfit@porttopm
  \global\textwidth=\elfit@porttw
  \global\textheight=\elfit@portth
  \global\hsize=\elfit@porttw
  \global\linewidth=\elfit@porttw
  \global\columnwidth=\elfit@porttw
  \global\vsize=\elfit@portth
  \global\@colht=\elfit@portth
  \global\@colroom=\elfit@portth
  \elfit@applypapersize
  \gdef\elfit@page{portrait}%
  \global\elfit@landscapedfalse
}

% The caption prints the declared type size. Fit never rewrites it
% down to the 7 pt floor.
\let\elfit@caption\sysml@caption
\renewcommand{\sysml@caption}{%
  \begingroup
    \ifx\elfit@shown\@empty\else
      \edef\sysml@body{\elfit@shown}%
    \fi
    \elfit@caption
  \endgroup
}

% Room for the sysml header, the 3 mm gap, and \topskip, so the header
% and the canvas together sit inside the chosen text height.
\def\elfit@reserve{%
  \setbox\elfit@picbox=\vbox{%
    \hsize=\textwidth
    \raggedright
    \noindent\sysml@caption
  }%
  \elfit@availht=\textheight
  \advance\elfit@availht by -\ht\elfit@picbox
  \advance\elfit@availht by -\dp\elfit@picbox
  \advance\elfit@availht by -3mm
  \advance\elfit@availht by -\topskip
  \advance\elfit@availht by -2pt
}

\long\def\elfit@measure#1{%
  \setbox\elfit@picbox=\hbox{%
    \elfit@tikzpicture[%
      x=\elfit@unit mm,%
      y=-\elfit@unit mm,%
      scale=\elfit@uniform,%
      transform shape%
    ]%
    \global\sysml@obscount=0\relax
    \ifsysml@marktrace
      \pgfpointtransformed{\pgfpointorigin}%
      \typeout{SYSMLORIGIN x=\the\pgf@x\space y=\the\pgf@y}%
    \fi
    #1%
    \elfit@endtikzpicture
  }%
}

\def\elfit@refuse#1{%
  \PackageError{endleaf}{#1}{Split the sheet. 7pt is a floor, not a fit target. Do not shrink type or the millimetre unit.}%
}

\def\elfit@fitcheck{%
  \elfit@nath=\dimexpr\ht\elfit@picbox+\dp\elfit@picbox\relax
  \ifdim\wd\elfit@picbox>\linewidth
    \elfit@refuse{picture does not fit \elfit@papername\space at the declared type size. Text was not shrunk}%
  \fi
  \ifdim\elfit@nath>\elfit@availht
    \elfit@refuse{picture does not fit \elfit@papername\space at the declared type size. Text was not shrunk}%
  \fi
}

\def\elfit@publishtype{%
  \xdef\elfit@shown{\elfit@bodyreq}%
}

\def\elfit@landreserve{%
  \setbox\elfit@picbox=\vbox{%
    \hsize=\elfit@landtw
    \raggedright
    \noindent\sysml@caption
  }%
  \elfit@availht=\elfit@landth
  \advance\elfit@availht by -\ht\elfit@picbox
  \advance\elfit@availht by -\dp\elfit@picbox
  \advance\elfit@availht by -3mm
  \advance\elfit@availht by -\topskip
  \advance\elfit@availht by -2pt
}

\long\def\elfit@fitcanvas#1{%
  \edef\elfit@bodyreq{\sysml@body}%
  \gdef\elfit@shown{\sysml@body}%
  \def\elfit@uniform{1}%
  \def\elfit@unit{1}%
  \ifsysml@marktrace \def\elfit@marksave{1}\else \def\elfit@marksave{0}\fi
  \sysml@marktracefalse
  \global\elfit@needlandfalse
  \elfit@reserve
  \elfit@measure{#1}%
  \ifdim\wd\elfit@picbox>\elfit@portline
    \ifelfit@locked
      \elfit@refuse{picture does not fit \elfit@pagename\space\elfit@papername\space at the declared type size. Text was not shrunk. Omit orientation to let Endleaf choose}%
    \else
      \global\elfit@needlandtrue
    \fi
  \fi
  \ifelfit@needland
    \gdef\elfit@page{landscape}%
    \ifelfit@insidefloat
      \global\elfit@thislandtrue
      \elfit@landreserve
      \elfit@measure{#1}%
      \elfit@nath=\dimexpr\ht\elfit@picbox+\dp\elfit@picbox\relax
      \ifdim\wd\elfit@picbox>\elfit@landline
        \elfit@refuse{picture does not fit landscape \elfit@papername\space at the declared type size. Text was not shrunk}%
      \fi
      \ifdim\elfit@nath>\elfit@availht
        \elfit@refuse{picture does not fit landscape \elfit@papername\space at the declared type size. Text was not shrunk}%
      \fi
    \else
      \elfit@setlandscape
      \elfit@reserve
      \elfit@measure{#1}%
      \elfit@fitcheck
    \fi
  \else
    \xdef\elfit@page{\elfit@pagename}%
    \elfit@nath=\dimexpr\ht\elfit@picbox+\dp\elfit@picbox\relax
    \ifdim\elfit@nath>\elfit@availht
      \elfit@refuse{picture does not fit \elfit@pagename\space\elfit@papername\space at the declared type size. Text was not shrunk}%
    \fi
  \fi
  \ifnum\elfit@marksave=1 \sysml@marktracetrue\fi
  \typeout{ENDLEAF_FIT_SHIP}%
  \elfit@measure{#1}%
  \elfit@publishtype
  \elfit@nath=\dimexpr\ht\elfit@picbox+\dp\elfit@picbox\relax
  \ifelfit@needland
    \edef\elfit@reportline{\the\elfit@landline}%
  \else
    \edef\elfit@reportline{\the\elfit@portline}%
  \fi
  \typeout{ENDLEAF_FIT shipped=\the\wd\elfit@picbox\space line=\elfit@reportline\space high=\the\elfit@nath\space avail=\the\elfit@availht\space uniform=\elfit@uniform\space unit=\elfit@unit\space type=\elfit@shown\space page=\elfit@page}%
  \ifdim\wd\elfit@picbox>\elfit@reportline\relax
    \typeout{ENDLEAF_OVERHANG package=sysml-tikz width=\the\wd\elfit@picbox\space line=\elfit@reportline}%
  \fi
  \leavevmode\box\elfit@picbox
}

\RenewDocumentEnvironment{sysmlcanvas}{+b}{%
  \elfit@fitcanvas{#1}%
}{}

% Header plus canvas stay together after the page-size decision, so
% the minipage uses the chosen linewidth. Do not open a portrait
% minipage before the canvas is measured.
\let\elfit@typeset\sysml@typeset
\renewcommand{\sysml@typeset}{%
  \noindent\begin{minipage}{\linewidth}%
    \elfit@typeset
  \end{minipage}%
}

% Do not collect the float body. \@float has already begun, so the
% minipage opens in internal vertical mode. A # in the body stays a #.
% A figure with no [..] (Pandoc's Figure) takes the class default. Expand
% \fps@figure first: \@xfloat reads the option list token by token.
\def\figure{\@ifnextchar[\elfit@figopt{\expandafter\elfit@figopt\expandafter[\fps@figure]}}
\def\elfit@figopt[#1]{%
  \global\elfit@insidefloattrue
  \global\elfit@thislandfalse
  \elfit@figure[#1]%
  \noindent\minipage{\linewidth}%
}
% \let keeps the source scan clean: a defined name must not start
% with end. The expansion is the same: close the minipage, then the float.
% ENDLEAF-R59-FLOAT. A landscape canvas marks its float before \end@float
% queues it: placement p only (24 = no "!" plus p), and the box is listed
% so the float-page builders below treat it as landscape.
\def\elfit@closefigure{%
  \endminipage
  \ifelfit@thisland
    \elfit@markland
  \fi
  \elfit@endfigure
  \global\elfit@thislandfalse
  \global\elfit@insidefloatfalse
}
\def\elfit@landname#1{elfit@land@\string#1}
\def\elfit@markland{%
  \@tempcnta\count\@currbox
  \divide\@tempcnta\@xxxii
  \multiply\@tempcnta\@xxxii
  \advance\@tempcnta 24\relax
  \global\count\@currbox\@tempcnta
  \expandafter\global\expandafter\let
    \csname\expandafter\elfit@landname\@currbox\endcsname\@empty
}
% A landscape box never joins a portrait float page: it fails the width
% test, so later figures wait behind it and the order is kept.
\let\elfit@testwrongwidth\@testwrongwidth
\def\@testwrongwidth#1{%
  \elfit@testwrongwidth#1%
  \expandafter\ifx\csname\elfit@landname#1\endcsname\@empty
    \global\@testtrue
  \fi
}
% When the head of the deferred list is a landscape box, the float page
% is that box alone, centred on the landscape text block. \@makefcolumn
% (\clearpage, \end{document}) calls this too.
\def\elfit@headof\@elt#1#2\@nil{\def\elfit@head{#1}}
\def\elfit@restof\@elt#1#2\@nil#3{\gdef#3{#2}}
\let\elfit@tryfcolumn\@tryfcolumn
\def\@tryfcolumn#1{%
  \global\@fcolmadefalse
  \ifx#1\@empty
  \else
    \expandafter\elfit@headof#1\@nil
    \expandafter\ifx\csname\expandafter\elfit@landname\elfit@head\endcsname\@empty
      \elfit@landcolumn#1%
    \else
      \elfit@tryfcolumn#1%
    \fi
  \fi
}
\def\elfit@landcolumn#1{%
  \global\@fcolmadetrue
  \global\setbox\@outputbox\vbox to\elfit@landth{%
    \vskip\@fptop
    \hbox to\elfit@landtw{\hfil\box\elfit@head\hfil}%
    \vskip\@fpbot}%
  \expandafter\elfit@restof#1\@nil#1%
  \begingroup
    \let\@elt\relax
    \xdef\@freelist{\@freelist\@elt\elfit@head}%
  \endgroup
  \expandafter\global\expandafter\let
    \csname\expandafter\elfit@landname\elfit@head\endcsname\relax
  \global\elfit@shiplandtrue
}
% The page that carries a landscape float page ships landscape; the
% geometry is portrait again before the next column starts.
\let\elfit@outputpage\@outputpage
\def\@outputpage{%
  \ifelfit@shipland
    \global\elfit@shiplandfalse
    \elfit@setlandscape
    \elfit@outputpage
    \elfit@setportrait
  \else
    \elfit@outputpage
  \fi
}
\let\endfigure\elfit@closefigure
\def\table{\@ifnextchar[\elfit@tabopt{\elfit@tabopt[\fps@table]}}
\def\elfit@tabopt[#1]{%
  \global\elfit@insidefloattrue
  \elfit@table[#1]%
  \noindent\minipage{\linewidth}%
}
\def\elfit@closetable{%
  \endminipage
  \elfit@endtable
  \global\elfit@insidefloatfalse
}
\let\endtable\elfit@closetable

% Plain pictures are not scaled. circuitikz is \let to \tikzpicture
% after that name is defined, so point it at this wrapper.
% ENDLEAF-R69-AUTO. With no job orientation, a plain picture wider than
% the portrait line at the declared type size turns its page landscape,
% as a sysml canvas does (ENDLEAF-R49, ENDLEAF-R56-FIT). Inside a figure
% the float is marked landscape (ENDLEAF-R59-FLOAT); elsewhere the page
% ships landscape and the next page is portrait again. Type stays at the
% declared size. A named orientation stays locked: no turn.
% ENDLEAF-R69-WIDE-WARN. A picture still wider than the line it is set on
% logs ENDLEAF_WIDE_PICTURE. The worker turns that into a layout_overhang
% job warning, which the gate reads as figure_wider_than_text.
\newdimen\elfit@plainline
\newdimen\elfit@plainht
\newif\ifelfit@plainstay
\ifdefined\Gscale@div\else\RequirePackage{graphicx}\fi
% ENDLEAF-R74-FIT. A plain picture wider than the portrait line that is also
% taller than a landscape text block cannot be turned: its landscape page
% would be blank and the picture would be clipped on the next page. Such a
% picture stays portrait and is scaled to the portrait line (and to the
% portrait text height), as long as the smallest text in it stays at or above
% 7 pt (ENDLEAF-R26); ENDLEAF_PICSCALE logs the scale. Below that floor it is
% set unscaled and logs ENDLEAF_WIDE_PICTURE. \ifelfit@plainstay is true when
% the picture does not turn its page: a named orientation, or too tall.
\def\elfit@plainscale{%
  \ifelfit@locked\elfit@plainstaytrue\else\elfit@plainstayfalse\fi
  \ifdim\wd\elfit@plainbox>\elfit@portline
    \ifelfit@locked
    \else
      \elfit@plainht=\dimexpr\ht\elfit@plainbox+\dp\elfit@plainbox\relax
      \ifdim\elfit@plainht>\elfit@landth
        \elfit@plainstaytrue
        \Gscale@div\elfit@ps{\elfit@portline}{\wd\elfit@plainbox}%
        \Gscale@div\elfit@pqh{\elfit@portth}{\elfit@plainht}%
        \ifdim\elfit@pqh\p@<\elfit@ps\p@\let\elfit@ps\elfit@pqh\fi
        \ifdefined\ellay@minfs\let\elfit@pfs\ellay@minfs\else\def\elfit@pfs{99}\fi
        \Gscale@div\elfit@pmin{7\p@}{\elfit@pfs\p@}%
        \ifdim\elfit@ps\p@<\elfit@pmin\p@
        \else
          \typeout{ENDLEAF_PICSCALE package=tikz width=\the\wd\elfit@plainbox\space line=\the\elfit@portline\space scale=\elfit@ps}%
          \setbox\elfit@plainbox\hbox to\elfit@portline{\hss\scalebox{\elfit@ps}{\box\elfit@plainbox}\hss}%
        \fi
      \fi
    \fi
  \fi
}
\def\tikzpicture{\@ifnextchar[{\elfit@tikzon}{\elfit@tikzon[]}}
\def\elfit@tikzon[#1]{%
  \setbox\elfit@plainbox=\hbox\bgroup
    \elfit@tikzpicture[#1]%
}
\def\elfit@plainfit{%
  \elfit@plainscale
  \elfit@plainline=\linewidth
  \ifdim\wd\elfit@plainbox>\elfit@portline
    \ifelfit@plainstay
    \else
      \ifelfit@insidefloat
        \global\elfit@thislandtrue
      \else
        \elfit@setlandscape
        \global\elfit@shiplandtrue
        % ENDLEAF-R69-WHERE: this page is landscape now; measure on its line.
        \elfit@plainline=\elfit@landline
      \fi
      \typeout{ENDLEAF_AUTO_LANDSCAPE package=tikz width=\the\wd\elfit@plainbox\space portline=\the\elfit@portline}%
    \fi
  \fi
  \ifelfit@thisland
    \elfit@plainline=\elfit@landline
  \fi
  \ifdim\wd\elfit@plainbox>\elfit@plainline
    \typeout{ENDLEAF_WIDE_PICTURE width=\the\wd\elfit@plainbox\space line=\the\elfit@plainline}%
  \fi
}
% ENDLEAF-R69-WHERE. Inside a landscape float the minipage is still the
% portrait line; a picture wider than it is centred on it, so the picture
% is centred on the landscape page and overhangs both sides evenly.
\def\elfit@plaincentre{%
  \ifelfit@thisland
    \ifdim\wd\elfit@plainbox>\linewidth
      \setbox\elfit@plainbox\hbox to\linewidth{\hss\box\elfit@plainbox\hss}%
    \fi
  \else
    \elfit@barecentre
  \fi
}
% ENDLEAF-R70-KINDS. A drawing kind (its preamble defines
% \EndleafBareCentre) centres a bare picture that turned its own page on
% the landscape line: it starts the line without an indent and overhangs
% both sides evenly when it is still wider. Without that flag (sysml,
% fulldoc) nothing here changes.
\ifdefined\EndleafBareCentre
  \def\elfit@barecentre{%
    \ifelfit@insidefloat\else
      \ifelfit@landscaped
        \ifdim\wd\elfit@plainbox>\elfit@portline
          \ifvmode\noindent\fi
          \setbox\elfit@plainbox\hbox to\linewidth{\hss\box\elfit@plainbox\hss}%
        \fi
      \fi
    \fi
  }
\else
  \let\elfit@barecentre\relax
\fi
\def\elfit@closetizpicture{%
    \elfit@endtikzpicture
  \egroup
  \elfit@plainfit
  \elfit@plaincentre
  % ENDLEAF-R74-FIT: no paragraph indent, which the box would add to its width.
  \ifvmode\noindent\fi\box\elfit@plainbox
}
\let\endtikzpicture\elfit@closetizpicture
\let\circuitikz\tikzpicture
\let\endcircuitikz\endtikzpicture

% ENDLEAF-R70-SHELL. A document-shell diagram fence reaches the page as an
% image in \pandocbounded; the lane templates point that at
% \elfit@bounded. The fence is drawn at the declared body size (10 pt
% unless \EndleafFenceTypePt says otherwise), so scaling by s sets its type
% at s times that size. 7 pt is the floor (ENDLEAF-R26).
%  1. It fits the portrait line: unchanged.
%  2. Shrinking to the portrait line keeps type at or above 7 pt: shrink,
%     and log the shrink.
%  3. Otherwise, with no job orientation: turn landscape (a float ships
%     alone on a landscape page; elsewhere this page ships landscape),
%     centre it, and shrink to the landscape line only while type stays at
%     or above 7 pt.
%  4. Otherwise (a named orientation, or still too wide landscape): set it
%     at the floor scale, centred; it overhangs and logs that.
% The height still bounds the scale, as Pandoc's own macro does.
% ENDLEAF-R80-CAPTION. The height bound B is \elfit@landth for a fence
% turned inside a float (that page is not set landscape yet, so
% \textheight is still the portrait value), else \textheight. Inside a
% wrapped figure or table float a caption is assumed, and the scale fits
% (ht+dp)*scale + \elfit@capreserve into B: the reserve is
% \ElfitCaptionHeight+\abovecaptionskip when the optional hook
% \ElfitCaptionHeight is defined, else 3\baselineskip (at \normalsize)
% +\abovecaptionskip. Outside a float the reserve is 0. The height scale
% never goes below the 7 pt floor; a fence still taller than B at the
% floor is set there and logs ENDLEAF_SHELLTALL (ENDLEAF-R80-TALL-WARN),
% height = (ht+dp)*scale + reserve, avail = B.
% Log lines (compile time, then shipout):
%   ENDLEAF_SHELLFIT n natural-width line type-pt scale
%   ENDLEAF_SHELLTALL n height avail
%   ENDLEAF_SHELLTURN n
%   ENDLEAF_SHELLPAGE n page portrait|landscape
\newcount\elfit@fencen
\newbox\elfit@fencebox
\newdimen\elfit@fenceline
\ifdefined\EndleafFenceTypePt
  \edef\elfit@fencetype{\EndleafFenceTypePt}
\else
  \def\elfit@fencetype{10}
\fi
\def\elfit@fencefloor{7}
\newdimen\elfit@fencebound
\newdimen\elfit@capreserve
\newdimen\elfit@tallht
% \elfit@fencescale{line}: \elfit@s = min(line / width, 1).
\def\elfit@fencescale#1{%
  \Gscale@div\elfit@s{#1}{\wd\elfit@fencebox}%
  \ifdim\elfit@s\p@>\p@\def\elfit@s{1}\fi
}
\def\elfit@bounded#1{%
  \global\advance\elfit@fencen\@ne
  \setbox\elfit@fencebox\hbox{#1}%
  \Gscale@div\elfit@minscale{\elfit@fencefloor\p@}{\elfit@fencetype\p@}%
  \elfit@fenceline=\linewidth
  \let\elfit@fenceturn\relax
  \elfit@fencescale\elfit@fenceline
  \ifdim\elfit@s\p@<\elfit@minscale\p@
    \ifelfit@locked
      \let\elfit@s\elfit@minscale
    \else
      \def\elfit@fenceturn{\typeout{ENDLEAF_SHELLTURN \the\elfit@fencen}}%
      \ifelfit@insidefloat
        \global\elfit@thislandtrue
      \else
        \ifelfit@landscaped\else
          \elfit@setlandscape
          \global\elfit@shiplandtrue
        \fi
      \fi
      \elfit@fenceline=\elfit@landline
      \elfit@fencescale\elfit@fenceline
      \ifdim\elfit@s\p@<\elfit@minscale\p@
        \let\elfit@s\elfit@minscale
      \fi
    \fi
  \fi
  \ifelfit@thisland
    \elfit@fenceline=\elfit@landline
  \fi
  % ENDLEAF-R80-CAPTION: height bound and caption reserve.
  \ifelfit@thisland
    \elfit@fencebound=\elfit@landth
  \else
    \elfit@fencebound=\textheight
  \fi
  \ifelfit@insidefloat
    \ifdefined\ElfitCaptionHeight
      \elfit@capreserve=\dimexpr\ElfitCaptionHeight+\abovecaptionskip\relax
    \else
      \begingroup
        \normalsize
        \xdef\elfit@captmp{\the\dimexpr 3\baselineskip+\abovecaptionskip\relax}%
      \endgroup
      \elfit@capreserve=\elfit@captmp
    \fi
  \else
    \elfit@capreserve=\z@
  \fi
  \let\elfit@fencetall\relax
  \Gscale@div\elfit@h{\dimexpr\elfit@fencebound-\elfit@capreserve\relax}{\dimexpr\ht\elfit@fencebox+\dp\elfit@fencebox\relax}%
  \ifdim\elfit@h\p@<\elfit@s\p@\let\elfit@s\elfit@h\fi
  \ifdim\elfit@s\p@<\elfit@minscale\p@
    \let\elfit@s\elfit@minscale
    \elfit@tallht=\elfit@s\dimexpr\ht\elfit@fencebox+\dp\elfit@fencebox\relax
    \advance\elfit@tallht\elfit@capreserve
    \ifdim\elfit@tallht>\elfit@fencebound
      \def\elfit@fencetall{\typeout{ENDLEAF_SHELLTALL \the\elfit@fencen\space\the\elfit@tallht\space\the\elfit@fencebound}}%
    \fi
  \fi
  \typeout{ENDLEAF_SHELLFIT \the\elfit@fencen\space\the\wd\elfit@fencebox\space\the\elfit@fenceline\space\elfit@fencetype\space\elfit@s}%
  \elfit@fencetall
  \elfit@fenceturn
  \ifdim\elfit@s\p@<\p@
    \setbox\elfit@fencebox\hbox{\scalebox{\elfit@s}{\box\elfit@fencebox}}%
  \fi
  \edef\elfit@tmp{\write16{ENDLEAF_SHELLPAGE \the\elfit@fencen\space\noexpand\thepage\space\noexpand\elfit@page}}%
  \ifx\elfit@fenceturn\relax
    % A shrunk fence fills the line: start it without the indent.
    \ifdim\elfit@s\p@<\p@\ifvmode\noindent\fi\fi
    \ifdim\wd\elfit@fencebox>\linewidth
      \setbox\elfit@fencebox\hbox to\linewidth{\hss\box\elfit@fencebox\hss}%
    \fi
  \else
    \ifvmode\noindent\fi
    \setbox\elfit@fencebox\hbox to\linewidth{\hss\box\elfit@fencebox\hss}%
  \fi
  \leavevmode\hbox{\elfit@tmp\box\elfit@fencebox}%
}
\makeatother
